% Section 6.3, from lectures/L06/README.md: what you should be able to do, the questions to test
% yourself with, and the next lecture.

\section{Sammanfattning}
\label{c6:sec:summary}

\needspace{6\baselineskip}
\subsection*{Det här ska du kunna nu}
Efter det här kapitlet ska du kunna:
\begin{itemize}
  \item Förklara skillnaden mellan ett kombinatoriskt nät och ett sekvensnät, och känna igen vilket
    en krets är.
  \item Förklara hur återkoppling ger minne, med självhållningen och SR-låset som exempel.
  \item Beskriva skillnaden mellan ett D-lås och en D-vippa, och rita utgången för båda i ett
    tidsdiagram.
  \item Räkna ut klockperioden från klockfrekvensen, och tvärtom.
  \item Förklara vad ett register, en räknare och ett skiftregister är, och vad de är byggda av.
  \item Rita mikrodatorns blockschema och redogöra för vad varje block gör: CPU, styrenhet,
    register, ALU, statusregister, programräknare, programminne, dataminne, bussar, portar och
    klocka.
  \item Beskriva instruktionscykeln och följa ett kort program instruktion för instruktion.
  \item Redogöra för ATmega328P:s minnen och portar, och vad som skiljer flash, SRAM och EEPROM
    åt.
\end{itemize}

\needspace{6\baselineskip}
\subsection*{Frågor att testa dig själv med}
\begin{enumerate}
  \item Varför går det inte att skriva en vanlig sanningstabell för ett SR-lås?
  \item Vad vinner man på att en vippa bara läser \code{D} vid den stigande flanken?
  \item Hur lång är en klockcykel vid 16~MHz?
  \item En räknare med 8 vippor står på sitt högsta värde och får en klockflank till. Vad händer?
  \item Vilket block i mikrodatorn räknar, och vilket bestämmer vad som ska räknas?
  \item Vad är ett hopp i ett program, sett från programräknarens håll?
  \item Vilket av ATmega328P:s minnen tappar sitt innehåll när strömmen bryts, och varför gör inte
    de andra det?
  \item Vad är en utport, sett som en byggsten av vippor?
\end{enumerate}

\needspace{6\baselineskip}
\subsection*{Nästa kapitel}
Kapitel~\ref{c7:ch} handlar om mikrodatorn på nära håll och om hur den programmeras:
\begin{itemize}
  \item ATmega328P på nära håll: de 32 registren, statusregistret och de tre minnena.
  \item Assemblerspråket: instruktioner, operander, direktiv och etiketter, och hur en rad blir
    maskinkod.
  \item Microchip Studio: skapa ett projekt, assemblera, och stega genom programmet i simulatorn.
  \item \textbf{Labb~3} börjar, bilaga~\ref{app:lab3}.
\end{itemize}
